\DocumentMetadata{lang=es-DO,testphase={phase-III,math}}
\documentclass[11pt]{scrartcl}
\usepackage{icv}

% -----------------------------------------------------------------------
% Metadata -- reemplaza estos valores para cada asignación nueva.
% -----------------------------------------------------------------------
\icvsetup{
  asignatura  = {ICV 442-T -- Acueductos y Alcantarillados},
  titulocorto = {Acueductos y Alcantarillados},
  titulo      = {Título de la asignación},
  subtitulo   = {Subtítulo o tema específico (opcional)},
  profesor    = {Dr. Ricardo Hernández Moreira},
  periodo     = {2026-1},
  version     = {v1},
}

\hypersetup{
  pdftitle    = {Título de la asignación},
  pdfauthor   = {Dr. Ricardo Hernández Moreira},
  pdfsubject  = {ICV 442-T -- Acueductos y Alcantarillados},
}

\begin{document}
\icvmaketitle

\begin{icvobjectives}
  \item Primer objetivo de aprendizaje que evalúa esta asignación.
  \item Segundo objetivo de aprendizaje.
\end{icvobjectives}

\section{Instrucciones}

Instrucciones generales: formato de entrega, fecha límite, trabajo
individual o en equipo, qué debe incluir el documento entregado.

\section{Problemas}

\begin{icvproblem}[points=15]{Cálculo de caudal de diseño}
Enunciado del primer problema. Puede incluir datos, tablas y figuras --
usa \texttt{icvincludegraphics} en vez de \texttt{includegraphics}, que
exige texto alternativo como tercer argumento:
\begin{verbatim}
\icvincludegraphics[width=4cm,alt={Descripción del contenido de la
  figura para lectores de pantalla.}]{ruta/figura.png}
\end{verbatim}
\end{icvproblem}

\begin{icvproblem}[points=10]{Verificación de presión}
Enunciado del segundo problema, sin figura.
\end{icvproblem}

\begin{icvproblem}{Pregunta de desarrollo (sin puntaje asignado)}
Enunciado del tercer problema -- \texttt{points} es opcional.
\end{icvproblem}

\end{document}
